 \documentclass[11pt,a4paper]{article}

\usepackage[top=2cm,bottom=2.5cm,left=2cm,right=2cm]{geometry}
\usepackage[dvipsnames,svgnames,x11names]{xcolor}
\usepackage{amsmath,amssymb}
\usepackage{tabularx}
\usepackage{colortbl}
\usepackage{tcolorbox}
\tcbuselibrary{skins}
\usepackage{hyperref}

\hypersetup{
    pdfauthor={Uniwersytet Vizja},
    pdftitle={Neuropsychologia - Ćwiczenie nr 1},
    hidelinks
}

\setlength{\parindent}{0pt}
\setlength{\parskip}{1em}

% Custom Colors
\definecolor{navyheader}{RGB}{19, 44, 89}
\definecolor{titleblue}{RGB}{23, 62, 115}
\definecolor{lightbgblue}{RGB}{243, 247, 253}
\definecolor{labelbgblue}{RGB}{218, 230, 246}
\definecolor{borderblue}{RGB}{160, 190, 225}

\begin{document}

% Top Meta / Running Header Context
\begin{center}
    {\footnotesize\textcolor{darkgray}{Neuropsychologia \textbar{} Ćwiczenie nr 1 \textbar{} Mapowanie mózg–zachowanie i analiza sieci neuronowych}}
    \vspace{0.2cm}
    \color{navyheader}\hrule height 1.5pt
\end{center}

\vspace{0.3cm}

% Header Block
\begin{tcolorbox}[
    enhanced,
    colback=lightbgblue,
    colframe=borderblue,
    arc=0mm,
    boxrule=1pt,
    left=15pt,
    right=15pt,
    top=0pt,
    bottom=15pt
]
    % Top solid banner
    \begin{tcolorbox}[
        enhanced,
        colback=navyheader,
        colframe=navyheader,
        arc=0mm,
        left=0pt,
        right=0pt,
        top=10pt,
        bottom=10pt,
        enlarge left by=-15pt,
        enlarge right by=-15pt,
        enlarge top by=0pt,
        width=\dimexpr\textwidth+30pt\relax
    ]
        \centering
        {\large\bfseries\textcolor{white}{UNIWERSYTET VIZJA}}\\[3pt]
        {\small\textcolor{white}{Wydział Psychologii i Nauk Społecznych}}
    \end{tcolorbox}

    \vspace{0.6cm}
    \centering
    {\LARGE\bfseries\textcolor{titleblue}{NEUROPSYCHOLOGIA}}\\[6pt]
    {\large\bfseries\textcolor{titleblue}{Ćwiczenie nr 1}}\\[8pt]
    {\normalsize Wprowadzenie do neuropsychologii klinicznej:}\\[4pt]
    {\large\bfseries\textcolor{titleblue}{Mapowanie relacji mózg–zachowanie i analiza sieci neuronowych}}
    \vspace{0.2cm}
\end{tcolorbox}

\vspace{0.5cm}

% Metadata Table
\renewcommand{\arraystretch}{1.4}
\noindent
\begin{tabularx}{\textwidth}{|>{\columncolor{labelbgblue}}p{3.5cm}|X<{\columncolor{lightbgblue}}|c@{}}
\hline
\textbf{Przedmiot:} & Neuropsychologia & \tabularnewline
\hline
\textbf{Rok akademicki:} & 2025/2026 & \tabularnewline
\hline
\textbf{Forma zajęć:} & Ćwiczenia & \tabularnewline
\hline
\textbf{Metody:} & Analiza przypadków klinicznych, praca grupowa, mapowanie neuroanatomiczne & \tabularnewline
\hline
\end{tabularx}

\vspace{0.6cm}

% Section 1
\section*{1. Informacje ogólne o zajęciach}

Pierwsze ćwiczenia laboratoryjne z neuropsychologii mają charakter wprowadzający, lecz zakładają zaawansowany poziom przygotowania teoretycznego właściwy dla roku IV studiów psychologicznych. Zamiast podejścia deskryptywnego, od studentów oczekuje się aktywnej krytycznej analizy: oceny siły dowodowej klasycznych przypadków klinicznych, konfrontacji modeli lokalizacyjnych z podejściem sieciowym (connectome-based models) oraz konstruowania hipotez neuropsychologicznych z uzasadnieniem metodologicznym. Praca opiera się na analizie czterech przypadków, mapowaniu neuroanatomicznym i strukturalnej dyskusji grupowej.

\vspace{0.2cm}

% Learning Objectives Box
\begin{tcolorbox}[
    enhanced,
    colback=lightbgblue,
    colframe=navyheader,
    arc=0mm,
    boxrule=1pt,
    title={\small\textbf{Cele kształcenia (zgodne z KRK -- poziom 6. EQF)}},
    coltitle=white,
    colbacktitle=navyheader,
    left=12pt,
    right=12pt,
    top=10pt,
    bottom=10pt
]
\normalsize
\textbf{Wiedza:} Student krytycznie porównuje paradygmaty lokalizacjonizmu, ekwipotencjalizmu i modeli sieci neuronowych; charakteryzuje teorie emocji i pamięci w kontekście neuropsychologicznym (Damasio, Squire, LeDoux); omawia metodologiczne ograniczenia badań opartych na jednostkowych przypadkach klinicznych.

\vspace{0.6em}

\textbf{Umiejętności:} Student konstruuje strukturalnie poprawne hipotezy neuropsychologiczne oparte na zasadzie podwójnej dysocjacji; mapuje uszkodzenia mózgu na schematy neuroanatomiczne z uzasadnieniem cytoarchitektonicznym (pola Brodmanna); ocenia, które komponenty funkcji poznawczych podlegają i nie podlegają modularnemu rozgraniczeniu.

\vspace{0.6em}

\textbf{Kompetencje społeczne:} Student identyfikuje etyczne dylematy towarzyszące wieloletniemu badaniu pacjentów z uszkodzeniami mózgu; prezentuje stanowisko analityczne w dyskusji plenarnej; wykazuje wrażliwość na podmiotowość pacjentów neurologicznych w narracji klinicznej.
\end{tcolorbox}

\subsubsection{Wymagania wstępne i zakładany poziom wiedzy}

Studenci przystępujący do ćwiczeń powinni biegle operować następującą wiedzą:

•Neuroanatomia funkcjonalna: budowa mózgowia (płaty korowe, struktury limbiczne, zwoje
podstawy, móżdżek), główne szlaki włókien (fasciculus arcuatus, corpus callosum, fornix),
cytoarchitektura Brodmanna (przynajmniej pola 4, 6, 9, 17, 22, 44, 45, 46).

•Neuroprzekaźnictwo: układy dopaminergiczny, noradrenergiczny, serotoninergiczny i cholinergiczny.
Ich projekte i rola w funkcjach poznawczych.

•Psychologia poznawcza: modele systemów pamięci (Squire, 1992; Tulving, 1972), modele uwagi
(Posner, Baddeley), koncepcje funkcji wykonawczych.

•Metodologia badań: logika wnioskowania przyczynowego, ograniczenia badań obserwacyjnych,
pojęcie trafności wewnętrznej i zewnętrznej.

\subsubsection{Harmonogram zajęć}

OrganizacyjnyRegulamin pracowni, program całoroczny, zasady zaliczenia, podział na
grupy analityczne

Wykład interaktywnyHistoria i paradygmaty neuropsychologii: od frenologii do connectomics;
modele lokalizacji, sieci i dynamiki; teoria markera somatycznego

Analiza przypadków
(I runda)Praca w czwórkach: analiza czterech kazusów (Gage, H.M., S.M.,
Tan/Broca) wg ustrukturyzowanego protokołu; wypełnianie arkuszy
analitycznych

Dyskusja plenarna (I)Prezentacja wniosków grupowych; krzyżowa konfrontacja modeli; pytania
prowadzącego wymagające uzasadnienia metodologicznego
Ćwiczenie
mapowaniaIndywidualne mapowanie uszkodzeń na schematy neuroanatomiczne + ćwiczenie komparatywne: analiza dysocjacji między przypadkami

Dyskusja plenarna
(II)Synteza: które domeny funkcjonalne podlegają modułowej organizacji? 

Ograniczenia metody kazusu; implikacje dla praktyki klinicznej

Podsumowani
eZadanie domowe, kryteria oceny, kwestie organizacyjne


\subsection{Podstawy teoretyczne}

\subsubsection{Czym jest neuropsychologia}
Neuropsychologia jest dyscypliną naukową i kliniczną leżącą na styku neurologii, psychiatrii i psychologii.

Przedmiotem jej badań są związki między strukturą i czynnościowym funkcjonowaniem mózgu a zachowaniem
– w tym procesami poznawczymi, emocjonalnymi i motywacyjnymi. W odróżnieniu od neurobiologii
kognitywnej, neuropsychologia kładzie szczególny nacisk na konsekwencje uszkodzeń mózgowych dla
codziennego funkcjonowania jednostki i dysponuje własnym zestawem metod diagnozy i rehabilitacji.
Jako dyscyplina kliniczna neuropsychologia obejmuje ocenę neuropsychologiczną (ang. neuropsychological
assessment), rehabilitację neuropsychologiczną i planowanie interwencji terapeutycznych. Wymaga zarówno
głębokiej znajomości anatomii i fizjologii układu nerwowego, jak i psychometrycznych podstaw diagnozy – co
czyni ją wyjątkowo wymagającą, lecz też interdyscyplinarnie bogatą dziedziną.
Neuropsychologia – Ćwiczenie nr 1 | Relacja mózg–zachowanie | Uniwersytet VizjaNeuropsychologia | Ćwiczenie nr 1 | Mapowanie mózg–zachowanie i analiza sieci neuronowych

\subsection{Rys historyczny i ewolucja paradygmatów}

\subsubsection{Frenologia – błąd z ważną intuicją}

Franz Joseph Gall (1758–1828) sformułował teorię frenologiczną, zakładającą korelację między topografią
czaszki a cechami psychicznymi. Choć frenologia okazała się pseudonauką, zasługuje na historyczną uwagę
jako pierwsza systematyczna próba przypisania funkcji psychicznych specyficznym obszarom mózgowia – co
stanowiło proto-lokalizacjonizm. Johann Caspar Spurzheim (1776–1832) popularyzował te idee w Europie
i Ameryce, co paradoksalnie przyczyniło się do stymulacji rzetelnych badań neuroanatomicznych jako
odpowiedzi na spekulacje Galla.

\subsubsection{Klasyczny lokalizacjonizm: Broca, Wernicke, Dejerine}

Paul Broca (1861) przedstawił przypadek Louisa Victora Leborgne'a (pacjent Tan) – mężczyzny, który po
uszkodzeniu lewego płata czołowego (okolica odpowiadająca polu 44 wg Brodmanna) utracił zdolność
artykulacji przy zachowanym rozumieniu mowy. Broca na podstawie dziewięciu kolejnych przypadków
sformułował tezę o lewostronnej dominacji językowej, stając się ojcem nowoczesnego lokalizacjonizmu.

Carl Wernicke (1874) opisał odmienny zespół: pacjenci z uszkodzeniem tylnego zakrętu skroniowego górnego
(pole 22) mówili płynnie, lecz ich mowa była pozbawiona sensu (parafazje, żargon). Wernicke zaproponował
pierwszy model połączeniowy (disconnection model): afazja przewodzenia wynikała jego zdaniem
z przerwania włókien łączących obie struktury językowe – czego wizualnym wyrazem jest fasciculus arcuatus.

Joseph Jules Dejerine (1891, 1892) uzupełnił ten obraz, opisując aleksję bez agrafii (uszkodzenie płata
potylicznego i ciała modzelowatego) oraz aleksję z agrafią (zakręt kątowy). Jego prace wyznaczyły kierunek
neuropsychologii połączeniowej, do której wrócimy omawiając model Geschwind-Wernicke.

\subsubsection{Holizm, ekwipotencjalizm i prawo akcji masowej}

Pierre Flourens (1794–1867) jako pierwszy podał w wątpliwość lokalizacjonizm na podstawie eksperymentów
ablacyjnych na ptakach: usunięcie różnych partii kory dawało podobne deficyty. Karl Lashley (1890–1958)
rozwinął tę intuicję w prawo ekwipotencjalizmu i prawo akcji masowej (mass action): zdolności uczenia się są
proporcjonalne do masy kory, a nie lokalizacji uszkodzenia. Chociaż prawo akcji masowej okazało się błędne
dla funkcji wysoce wyspecjalizowanych, pozostaje ważnym ostrzeżeniem przed naiwnym lokalizacjonizmem.

\subsubsection{Neuropsychologia sieciowa i era connectomics}
Współczesna neuropsychologia – wsparta neuroobrazowaniem czynnościowym (fMRI, PET) i strukturalnym
(DTI traktografia, VBM) – zajmuje stanowisko sieciowe: poszczególne funkcje poznawcze i emocjonalne
wymagają skoordynowanej aktywności rozproszonych sieci neuronowych. Wyróżniamy trzy kluczowe sieci
wielkich skal (large-scale networks):

•Sieć spoczynkowa (Default Mode Network, DMN): przyśrodkowa kora przedczołowa (mPFC), tylna
kora obręczy (PCC), zakręt nadbrzeżny, hipokamp. Aktywna w spoczynku (ruminacje, projekcja
w przyszłość), deaktywowana podczas zadań wymagających skupionej uwagi. Zaangażowana
w pamięć autobiograficzną, teorię umysłu i poczucie własnej tożsamości.

•Sieć salience (SN / Sieć istotności): przednia wyspa (anterior insula), przednia część zakrętu obręczy
(ACC, pole 32), ciało migdałowate. Odpowiada za detekcję bodźców biologicznie istotnych,
przełączanie między sieciami (DMN vs CEN) i interocepcję.

•Sieć kontroli centralnej (Central Executive Network, CEN): grzbietowo-boczna kora przedczołowa
(dlPFC), tylna kora ciemieniowa (PPC). Zaangażowana w planowanie, pamięć roboczą i regulację
odgórną.

Dynamiczne wzajemne hamowanie DMN i CEN – regulowane przez SN – stanowi neuronalne podłoże
przełączania między przetwarzaniem wewnętrznym a ukierunkowanym na zadanie. Zakłócenie tych relacji
obserwujemy w schizofrenii, ADHD i depresji.

\subsection{Kluczowe koncepcje neuropsychologiczne}

\subsubsection{Plastyczność neuronalna – mechanizmy i implikacje kliniczne}

Plastyczność neuronalna to zdolność układu nerwowego do reorganizacji w odpowiedzi na doświadczenia,
uczenie się lub uszkodzenie. Wyróżniamy hierarchię poziomów plastyczności:

•Plastyczność synaptyczna: długotrwałe wzmocnienie (LTP) i długotrwałe osłabienie (LTD) na poziomie
synaps glutamatergicznych (receptory AMPA/NMDA). Podstawowy mechanizm uczenia się i pamięci.

•Plastyczność strukturalna: zmiany w morfologii dendrytów (kolce dendrytyczne), tworzenie nowych
połączeń aksonalnych (sprouting) oraz reorganizacja map korowych po deprywacji sensorycznej lub
uszkodzeniu (przesunięcie mapy korowej – cortical remapping).

•Neurogeneaza dorosłych: powstawanie nowych neuronów w zakręcie zębatym hipokampu (granule
cells) i opuszce węchowej. Regulowana przez wysiłek fizyczny, stres i czynnik BDNF; hamowana przez
glikokortykoidy.

Kliniczne znaczenie plastyczności jest dwustronne: stanowi ona podstawę rehabilitacji neurologicznej (constraint-induced movement therapy, terapia mowy), lecz również podłoże maladaptacyjnych zmian – takich jak ból fantomowy, uzależnienia czy epileptogeneza.

\subsubsection{Diaschiza i ewolucja objawów po uszkodzeniu mózgu}

Diaschiza (gr. diaschizein – rozcinać) to zjawisko nagłego, przejściowego zaniku czynności w obszarach anatomicznie odległych od ogniska uszkodzenia, lecz z nim czynnościowo połączonych. Constantin von Monakow (1914) wprowadził to pojęcie, obserwując, że po ogniskowych uszkodzeniach korowych następuje czasowa utrata aktywności w obszarach podkorowych połączonych szlakami projekcyjnymi.

Diaschiza krzyżowa móżdżku (crossed cerebellar diaschisis) jest najlepiej udokumentowanym jej wariantem: uszkodzenie jednej półkuli mózgu powoduje zmniejszenie przepływu krwi i metabolizmu w przeciwległej półkuli móżdżku, ze względu na przekrzyżowanie szlaków korowo-mostowo-móżdżkowych. Zjawisko to ma bezpośrednie implikacje prognostyczne: wczesna poprawa po udarze może być wynikiem cofania się diaschizji, a nie trwałej regeneracji tkankowej.

\subsubsection{Podwójna dysocjacja – logika i ograniczenia}

Podwójna dysocjacja (double dissociation) stanowi fundament wnioskowania neuropsychologicznego o modularnej organizacji umysłu. Jeżeli pacjent A wykazuje deficyt funkcji X przy zachowanej Y, a pacjent B - deficyt Y przy zachowanej X, wnioskujemy, że mechanizmy neurobiologiczne X i Y są od siebie niezależne. Metoda ta jest argumentem silniejszym niż sama dysocjacja jednostronna, ponieważ wyklucza wyjaśnienie różnicą w trudności zadań.

Ograniczenia metodologiczne podwójnej dysocjacji są jednak poważne: (1) Uszkodzenia mózgu rzadko są
precyzyjnie ogniskowe – masa białej istoty jest zwykle uszkodzona równolegle, co zaciera granice modułów.
(2) Kompensacja funkcjonalna może maskować rzeczywiste deficyty. (3) Wnioskowanie o normalnej
architekturze mózgu na podstawie patologii jest logicznie problematyczne (argument Poeppela). (4) Próby n=1
uniemożliwiają generalizację statystyczną.


\subsubsection{Teoria markera somatycznego (Damasio, 1994)}
Antonio Damasio zaproponował teorię markera somatycznego jako wyjaśnienie mechanizmu podejmowania
decyzji w warunkach niepewności. Zgodnie z nią, stany somatyczne (zmiany fizjologiczne ciała – tętno,
przewodnictwo skórne, napięcie mięśniowe) są przywoływane podczas oceny opcji decyzyjnych i działają jako
szybkie, przedświadome sygnały wartościowania – markery. Strukturami kluczowymi dla tego mechanizmu są
brzuszno-przyśrodkowa kora przedczołowa (vmPFC) i ciało migdałowate (amygdala), łączące reprezentacje
somatyczne z wartością emocjonalną opcji.

Empirycznym dowodem na teorię jest Iowa Gambling Task (IGT; Bechara i in., 1994): pacjenci z uszkodzeniami
vmPFC konsekwentnie wybierają talie z natychmiastową nagrodą i dużą długoterminową stratą, pomimo
świadomości niekorzystności wyboru – brakuje im anticipatory skin conductance response, sygnalizującego
ryzyko. Przypadek Gage'a jest klasyczną ilustracją kliniczną: niezdolność do integracji emocjonalnych sygnałów
wartościujących skutkuje irracjonalnymi decyzjami życiowymi przy zachowanej inteligencji formalnej.

\subsection{Metoda przypadku klinicznego (n=1) – możliwości i granice}
Klasyczna neuropsychologia kliniczna w znacznej mierze opiera się na dogłębnej analizie jednostkowych
przypadków. Metoda ta ma zarówno unikatowe zalety, jak i poważne ograniczenia, które student roku IV musi
umieć ocenić:

Kryterium
Mocne strony metody n=1Ograniczenia metody n=1

Siła dowodowa
Możliwość szczegółowej, wieloletniej obserwacji; wykrycie rzadkich dysocjacji niemożliwych w próbach grupowychBrak możliwości generalizacji; ryzyko konfundowania przez cechy indywidualne pacjenta

Kontrola zmiennych
Dokładna charakterystyka deficytu przy znanych warunkach uszkodzeniaUszkodzenia in vivo rzadko są precyzyjnie ogniskowe; lezja może obejmować istotę białą

Implikacje teoretyczne
Jeden dobrze udokumentowany przypadek może obalić teorięWnioskowanie o normalnym mózgu z patologii jest logicznie ryzykowne (argument Poeppela)

Kontekst historyczny
Wiele fundamentalnych odkryć (Broca,Wernicke, H.M.) pochodzi z n=1
Możliwość błędu selekcji: opisuje się przypadki potwierdzające teorię, pomija anomalie

Etyka
Możliwe dogłębne uzyskanie zgody i budowanie relacji terapeutycznejRyzyko uprzedmiotowienia pacjenta przez jego długotrwałe badanie; kwestia anonimowości


\section{Pnaliza przypadków klinicznych}
Poniżej przedstawiono cztery przypadki kliniczne o historycznym i teoretycznym znaczeniu dla
neuropsychologii. W grupach czteroosobowych analizujecie każdy kazus, wypełniając protokół analityczny
(sekcja 5). Uwaga: opisy celowo zawierają zarówno fakty dobrze udokumentowane, jak i elementy historycznie
dyskusyjne – Waszym zadaniem jest ich identyfikacja i krytyczna ocena.

\subsubsection*{Przypadek 1 – Phineas Gage (1823–1860): Osobowość i kora przedczołowa}
Tło: W 1848 roku Phineas Gage, 25-letni brygadzista przy budowie kolei w Vermont, doznał urazu penetrującego:
metalowy pręt (109 cm, 6 kg) przebił jego czaszkę. Gage przeżył bez widocznych deficytów neurologicznych (ruch,
mowa, pamięć, percepcja pozostały nienaruszone), lecz jego osobowość i zachowanie uległy radykalnej zmianie.

Obraz kliniczny: Przed wypadkiem – sumienny, odpowiedzialny, szanowany. Po wypadku – nieodpowiedzialny,
impulsywny, wulgarny, niezdolny do planowania i realizacji celów; Harlow (1868) opisał go jako kogoś, dla kogo
"równowaga między zdolnościami intelektualnymi a zwierzęcą namiętnością" została zakłócona.

Neuroanatomia: Rekonstrukcja komputerowa czaszki Gage'a (Damasio i in., 1994; Van Horn i in., 2012) wskazuje na
uszkodzenie lewego i prawego płata czołowego – przede wszystkim vmPFC i OFC (pola Brodmanna 10, 11, 12, 25)
oraz połączeń z przednim zakrętem obręczy (ACC).

Kontrowersje: Późniejsze analizy historyczne (Macmillan, 2000) sugerują, że relacje o degeneracji osobowości Gage'a
mogły być przesadzone lub ideologicznie motywowane. Gage podobno odzyskał część funkcji społecznych i pracował
jako woźnica w Chile. Implikacja metodologiczna: nawet klasyczne przypadki podlegają rewizji.

\subsubsection*{Przypadek 2 – Henry Molaison (H.M., 1926–2008): Pamięć i hipokamp}
Tło: H.M. cierpiał na lekooporną padaczkę skroniową od 16. r.ż. W 1953 roku neurochirurg William Beecher Scoville
wykonał obustronną medialną resekcję skroniową, usuwając ok. 8 cm tkanki: hipokamp, korę śródwęchową
(entorhinal cortex) i parasubiculum, ciało migdałowate oraz korę śródskroniową (perirhinal cortex).

Obraz kliniczny: Natychmiastowa i kompletna amnezja następcza (anterograde amnesia), czyli niemożność tworzenia
nowych wspomnień deklaratywnych (epizodycznych i semantycznych). Amnezja wsteczna obejmowała ok. 11 lat
przed operacją (gradient temporalny). Pamięć proceduralna (mirror tracing, rotary pursuit) pozostała nienaruszona.
H.M. uczył się nowych umiejętności ruchowych bez świadomości tego faktu.

Znaczenie teoretyczne: Przypadek H.M. dostarczył kluczowych dowodów na dysocjację pamięci deklaratywnej (hippocampus-dependent) i proceduralnej (striatum/cerebellum-dependent), dał podstawy taksonomii Squire'a (1992) i wpłynął na rozumienie konsolidacji pamięci (Frankland & Bontempi, 2005).

Implikacja sieci: Uszkodzenie hipokampu zaburza nie pojedynczy moduł, lecz węzeł krytyczny w sieci pamięci obejmującej korę śródwęchową, korę asocjacyjną i ciało migdałowate. Amnezja H.M. była selektywna nie dlatego, że pamięć jest "w hipokampie", lecz dlatego, że hipokamp jest niezbędnym interfejsem między pamięcią roboczą,  a kodowaniem długotrwałym.

\subsubsection*{Przypadek 3 – S.M. (anonimowa pacjentka): Strach i ciało migdałowate}
Tło: S.M. jest kobietą z chorobą Urbacha-Wiethe (lipoidoza skórno-śluzówkowa, OMIM #247100) – autosomalnie
recesywnym schorzeniem powodującym stopniowe zwapnienie i destrukcję ciała migdałowatego (amygdala).

Przypadek badany przez Adolphsa, Damasio i współpracowników od lat 90. XX w.

Obraz kliniczny: S.M. nie doświadcza subiektywnego strachu nawet w sytuacjach obiektywnie zagrażających życiu
(agresywne węże i pająki, sceny z horrorów, napad z bronią, COVID-19). Rozpoznawanie emocji z twarzy jest
wybiórczo zaburzone dla strachu; inne emocje podstawowe pozostają relatywnie niezmienione. IQ, pamięć werbalna
i funkcje wykonawcze są w normie.

Paradoks: W 2021 roku Feinstein i in. wykazali, że S.M. może doświadczać paniki wywołanej inhalacją CO₂ – bodźcem
działającym przez interoceptywne szlaki wewnętrzne (jądro migdałowate łożyska i pień mózgu), a nie przez zmysły
zewnętrzne. Sugeruje to, że ciało migdałowate nie jest jedyną strukturą strachu, lecz kluczowym węzłem strachu
eksteroceptywnego.

Teoria LeDoux vs Damasio: Przypadek S.M. jest często cytowany jako potwierdzenie modelu LeDoux'a (amygdala
jako centrum strachu), lecz wymaga reinterpretacji w świetle danych o CO₂. Stanowi znakomity przykład przypadku,
który był zbyt pochopnie wyjaśniany przez istniejące teorie.

\subsubsection*{Przypadek 4 – Louis Victor Leborgne "Tan" (1809–1861): Język i asymetria półkulowa}
Tło: Leborgne trafił do szpitala Bicêtre w 1840 r. z powodu postępującej utraty mowy. Przez 21 lat potrafił
wypowiedzieć jedynie sylabę "tan" (stąd pseudonim) oraz przekleństwo – przy pełnym rozumieniu mowy
i sprawności niewerbalnej. Prawostronny niedowład kończyn sugerował ogniskowe uszkodzenie lewej półkuli. Paul
Broca zbadał go na łożu śmierci i wykonał sekcję mózgu w dniu po zgonie.

Neuroanatomia: Broca zlokalizował uszkodzenie w tylnej części trzeciego zakrętu czołowego lewej półkuli (gyrus
frontalis inferior, pars opercularis) – dziś znane jako pole Broki (pole Brodmanna 44). Późniejsze analizy MRI czaszki

Tan (Dronkers i in., 2007) wykazały jednak rozleglejsze uszkodzenie obejmujące wyspę i podkorową istotę białą, co
komplikuje czystą lokalizację.

Znaczenie historyczne: Przypadek Tan jest symbolem narodzin lokalizacjonizmu neurologicznego. Jednocześnie jest
doskonałym przykładem tego, jak posthumna analiza podważa proste interpretacje – uszkodzenie Tan było
rozleglejsze niż sugerował Broca, a afazja niefluentna jest funkcją wielu ośrodków i połączeń, nie jednego punktu.

Pytanie porównawcze: Zestawiając przypadek Tan z H.M. – w obu doszło do głębokiego i selektywnego deficytu
domeny poznawczej (język vs pamięć deklaratywna) przy zachowaniu innych funkcji. Czy stanowią one dowód na
modularną organizację umysłu? Jakie dodatkowe dane byłyby potrzebne, by odpowiedzieć twierdząco?


\section{Ćwiczenia: Mapowanie i analiza komparatywna}

\subsection{Ćwiczenie A – Mapowanie neuroanatomiczne (indywidualne)}

Na podstawie schematów bocznego i przyśrodkowego widoku mózgu (Załącznik A) wykonaj następujące
zadania:

1. Zaznacz i opisz na schemach obszary uszkodzone u każdego z czterech pacjentów (Gage, H.M., S.M.,
Tan), stosując odrębne oznaczenia literowe lub kolorystyczne. Dla każdego obszaru podaj: nazwę
polską, nazwę angielską, numer pola Brodmanna (jeśli dotyczy).

2. Nanieś na schemat trzy wielkie sieci (DMN, SN, CEN) z ich głównymi węzłami. Zaznacz, które węzły
sieci są uszkodzone w poszczególnych przypadkach.

3. Dla każdego przypadku sformułuj hipotezę neuropsychologiczną w strukturze: "Jeżeli [uszkodzona
struktura / węzeł sieci] odpowiada za [funkcję], to oczekujemy deficytu [X] przy zachowanej funkcji
[Y]". Hipoteza musi być falsyfikowalna.

4. Oceń dla każdego przypadku: (a) na ile uszkodzenie jest ogniskowe i daje się wyjaśnić modelem
lokalizacjonistycznym; (b) jakie aspekty obrazu klinicznego wymagają wyjaśnienia przez model
sieciowy. Podaj uzasadnienie neuroanatomiczne.


\subsection{Ćwiczenie B – Analiza dysocjacji komparatywna (grupowe)}
W grupach czteroosobowych wypełnijcie poniższą macierz dysocjacji. Dla każdej pary przypadków oceńcie, czy obserwujemy dysocjację prostą, podwójną lub brak dysocjacji – i uzasadnijcie ocenę.
Domena funkcjonalna
Gage vmPFC/OFC
H.M. Hipokamp
S.M. Amygdala
Tan Pole 44
Pamięć deklaratywna (epizodyczna)
Pamięć proceduralna / uczenie reguł
Planowanie i funkcje wykonawcze
Rozpoznawanie i przeżywanie strachu
Produkcja mowy (fluencja)
Rozumienie mowy
Regulacja decyzji opartych na emocjach
Tożsamość / osobowość
Legenda: Z = zachowana; D = deficyt; N = nieznana / niezdiagnozowana; (~) = częściowo zachowana lub zmienna

\subsection{Kryteria oceny ćwiczeń}
Kryterium
Poprawność anatomiczna
mapowania (ćw. A)
Waga (%)

20. Opis oczekiwanej jakości 

Precyzja lokalizacji z uwzględnieniem pól Brodmanna; poprawne oznaczenie węzłów sieci

Kryterium

Waga (%)

Opis oczekiwanej jakości

Jakość hipotezy neuropsychologicznej (ćw. A)25Falsyfikowalność, precyzja predykcji, poprawność logiczna struktury hipotezy

Ocena modeli: lokalizacjonizm vs sieciowy (ćw. A)20Zrównoważona, uzasadniona neuroanatomicznie ocena obu paradygmatów

Macierz dysocjacji – trafność i uzasadnienie (ćw. B)20Poprawna identyfikacja dysocjacji prostych i podwójnych; jakość uzasadnienia

Identyfikacja ograniczeń metodologicznych (ćw. B)15Krytyczna ocena siły dowodowej kazusów; wskazanie alternatywnych wyjaśnień

\section{Protokół analityczny studenta}

Wypełnij poniższy protokół w trakcie i po zakończeniu ćwiczeń. Odpowiedzi powinny mieć charakter analityczny – odwołuj się do konkretnych struktur, sieci i teorii, cytuj autorów (nazwisko, rok). Protokół stanowi element zaliczenia zajęć.

Imię i nazwisko

Numer albumu

Ocena końcowa

\subsection{Analiza przypadku 1 – Phineas Gage}

1. Wymień uszkodzone struktury (nazwy polskie, angielskie, pola Brodmanna) i przyporządkuj je do węzłów jednej z trzech wielkich sieci (DMN/SN/CEN).

2. Sformułuj hipotezę neuropsychologiczną wg wzoru: "Jeżeli [struktura] odpowiada za [funkcję], to oczekujemy [deficyt] przy zachowanej [funkcja Y]". Oceń falsyfikowalność swojej hipotezy.

3. Jak teoria markera somatycznego Damasio wyjaśnia zmiany zachowania Gage'a lepiej niż prosty lokalizacjonizm? Co w tym wyjaśnieniu pozostaje spekulatywne?

4. Macmillan (2000) sugeruje, że relacje o Gage'u były ideologicznie zniekształcone. Jakie implikacje ma to dla siły dowodowej tego przypadku jako argumentu w neuropsychologii?

5. Który komponent teorii sieci (DMN/SN/CEN) najbardziej ucierpiał w przypadku Gage'a i jak to wyjaśnia obserwowane objawy behawioralne?

\subsection{Analiza przypadku 2 – Henry Molaison (H.M.)}

1. Wymień wszystkie usunięte struktury i opisz ich funkcje wg
aktualnej wiedzy. Która z usuniętych struktur jest Twoim zdaniem
krytyczna dla obserwowanego deficytu i dlaczego?

2. Wyjaśnij mechanizm amnezji następczej: dlaczego H.M. mógł
przechowywać wspomnienia krótkoterminowe, lecz nie
konsolidować ich w długotrwałych? Odwołaj się do modelu
Squire'a (1992).

3. Dlaczego pamięć proceduralna H.M. pozostała nienaruszona?
Jakie struktury i szlaki ją podtrzymują? Jak to samo zjawisko
wyjaśniają modele sieciowe (np. sieć prążkowia)?

4. Dronkers i in. (2007) ponownie przeanalizowali przypadek Tan
MRI pośmiertnym – analogicznie Corkin powtórnie analizowała
H.M. Jakie ryzyko metodologiczne wiąże się z wielokrotną
reinterpretacją danych z jednego przypadku?

5. Sformułuj pytanie badawcze, które mogłoby uzupełnić wiedzę
uzyskaną z przypadku H.M. – i zaproponuj metodę (badanie
grupowe, fMRI, optogenetyka na modelu zwierzęcym), która
mogłaby je zbadać.

\subsection{Analiza przypadku 3 – S.M.}


1. Opisz mechanizm uszkodzenia ciała migdałowatego w chorobie
Urbacha-Wiethe. Co różni ten typ uszkodzenia od urazu lub
udaru? Jaki ma to wpływ na interpretację wyników?

2. Feinstein i in. (2021) wykazali, że S.M. doznaje paniki po inhalacji
CO₂. Jak ten wynik zmienia Twoją ocenę przypadku S.M. jako
dowodu na teorię LeDoux'a (amygdala = centrum strachu)?

3. Jak zinterpretować fakt, że S.M. rozumie pojęcie strachu
poznawczo i może opisać sytuacje strachogenne, lecz nie
doświadcza strachu afektywnie? Odwołaj się do debaty James-
Lange vs Cannon-Bard.

4. Powiąż uszkodzenie ciała migdałowatego z zaburzeniami sieci SN
(Sieć istotności). Jak zakłócenie SN może wyjaśniać trudności S.M.
w detekcji biologicznie istotnych zagrożeń?

5. Wskaż trzy etyczne dylematy związane z kilkudziesięcioletnim
badaniem pacjentki S.M. Jak można by minimalizować szkodę
etyczną przy jednoczesnym zachowaniu wartości naukowej
badań?

\subsection{Analiza przypadku 4 – Louis Victor Leborgne (Tan)}

1. Porównaj lokalizację uszkodzenia Tan wg Broki (1861) i wg Dronkers i in. (2007). Co zmieniają te nowe dane dla klasycznej interpretacji przypadku?

2. Afazja Broki jest niefluentna, lecz ze względnie zachowanym rozumieniem. Afazja Wernickego – odwrotnie. Jaką podwójną dysocjację stanowi para Tan / pacjent Wernickego? Jakie wnioski o architekturze systemu językowego z niej wynikają?



3. Model Geschwind-Wernicke zakłada kluczową rolę fasciculus arcuatus w transmisji informacji między obszarami Broki i Wernickego. Jak ten model wyjaśnia afazję przewodzenia? Gdzie są jego ograniczenia w świetle współczesnej traktografii DTI?

4. Porównaj przypadek Tan z przypadkiem H.M.: oba dotyczyły dramatycznych deficytów domenowo-specyficznych. Na ile oba przypadki łącznie stanowią argument za modularną organizacją umysłu? Jakie kontrargumenty możesz sformułować?

\subsection{Synteza komparatywna i dyskusja plenarna}

1. Na podstawie macierzy dysocjacji (ćw. B): wskaż wszystkie pary przypadków, które tworzą podwójną dysocjację. Dla każdej pary sformułuj wniosek o niezależności dwóch domen funkcjonalnych.

2. Który z czterech przypadków Twoim zdaniem najsłabiej nadaje się do wnioskowania o normalnej architekturze mózgu? Uzasadnij, odwołując się do kryteriów metodologicznych (ogniskowość uszkodzenia, kontrola konfundantów, gradient amnezji itp.).

3. Czy model sieciowy (DMN/SN/CEN) lepiej niż lokalizacjonizm wyjaśnia złożoność objawów w którymkolwiek z czterech przypadków? Podaj konkretny przykład z uzasadnieniem neuroanatomicznym.

4. Sformułuj syntezę: co cztery omówione przypadki łącznie mówią nam o relacji mózg–zachowanie? Jakie pytania pozostają nierozstrzygnięte i dlaczego są trudne do zbadania metodą n=1?

5. Zaproponuj jedno nowoczesne badanie neuroobrazowe (fMRI, DTI, MEG lub inne), które uzupełniłoby wiedzę o mechanizmach zaangażowanych w jeden z czterech przypadków. Opisz hipotezę, metodę i oczekiwany wynik.

\section{Zadanie domowe i literatura}

\subsection*{Zadanie domowe}

Napisz strukturalne opracowanie (min. 800 słów, APA 7. ed.) jednego z poniższych przypadków klinicznych:
(a) Zespół Capgrasa,
(b) Somatoparafrenia,
(c) Syndrom obcej ręki (alien hand syndrome),
(d) Prozopagnozja (przypadek P.S. lub inny udokumentowany).

Opracowanie musi zawierać:
(1) opis objawów klinicznych z terminologią specjalistyczną;
(2) proponowaną lokalizację uszkodzeń i zaangażowane sieci (DMN/SN/CEN jeśli dotyczy);
(3) krytyczną ocenę adekwatności modelu lokalizacjonistycznego vs sieciowego;
(4) omówienie siły dowodowej dostępnych danych (n=1 vs badania grupowe);
(5) refleksję nad etycznymi aspektami badania tego schorzenia.

Forma: .pdf lub .docx, 12 pt Calibri, marginesy 2,5 cm, przypisy w formacie APA 7.

\subsection{Literatura obowiązkowa}

1.Kolb, B. i Whishaw, I. Q. (2021). Fundamentals of Human Neuropsychology (8th ed.). Worth Publishers.
[Rozdziały 1–4, 16]

2.Damasio, A. R. (1994). Descartes' Error: Emotion, Reason and the Human Brain. Putnam. [Rozdziały 1–4]

3.Luria, A. R. (1976). The Working Brain: An Introduction to Neuropsychology. Basic Books. [Rozdziały 1–2]

4.Squire, L. R. (1992). Memory and the hippocampus: A synthesis from findings with rats, monkeys, and humans. Psychological Review, 99(2), 195–231.

\subsection{Literatura uzupełniająca – przypadki i metodologia}

Damasio, H., Grabowski, T., Frank, R., Galaburda, A. M. i Damasio, A. R. (1994). The return of Phineas Gage: Clues about the brain from the skull of a famous patient. Science, 264(5162), 1102–1105.

Macmillan, M. (2000). An Odd Kind of Fame: Stories of Phineas Gage. MIT Press.

Corkin, S. (2002). What's new with the amnesic patient H.M.? Nature Reviews Neuroscience, 3, 153–160.

Dronkers, N. F., Plaisant, O., Iba-Zizen, M. T. i Cabanis, E. A. (2007). Paul Broca's historic cases: High resolution MR imaging of the brains of Leborgne and Lelong. Brain, 130(5), 1432–1441.

Feinstein, J. S. i in. (2021). Fear and panic in humans with bilateral amygdala damage. Nature Neuroscience, 24,1. [Badanie CO₂ u S.M.]

Adolphs, R., Tranel, D., Damasio, H. i Damasio, A. (1994). Impaired recognition of emotion in facial expressions following bilateral damage to the human amygdala. Nature, 372, 669–672.


\subsection{Literatura uzupełniająca – modele sieciowe i metodologia}

Buckner, R. L., Andrews-Hanna, J. R. i Schacter, D. L. (2008). The brain's default network: Anatomy, function, and relevance to disease. Annals of the New York Academy of Sciences, 1124, 1–38.

Menon, V. (2011). Large-scale brain networks and psychopathology: A unifying triple network model. Trends in Cognitive Sciences, 15(10), 483–506.

Sporns, O. (2010). Networks of the Brain. MIT Press.

Poeppel, D. (1996). A critical review of PET studies of phonological processing. Brain and Language, 55(3), 317-385. [Argument metodologiczny]


\end{document}
